\chapter{Results}
\label{ch:results}
% ============================================================

\todo{Present quantitative and qualitative results objectively. Use tables and
figures. Do not interpret here — save that for Chapter 6.}

\section{Efficiency Metrics}

\subsection{Task Completion Time}
\subsection{Perceived Text Quality}
Task Completion Time
Dein Ansatz ist richtig. Ein paar Ergänzungen:
Normalisierung ist wichtig
Rohe Zeit allein ist trügerisch — ein langer Artikel dauert immer länger. Du solltest also normalisieren, z.B.:

Zeit pro 100 Wörter Output
Zeit pro Komplexitätsstufe (falls du Aufgaben vordefinierst)

Was du kontrollieren musst

Textlänge des Outputs
Themen-/Aufgabenschwierigkeit (möglichst vergleichbare Tasks in Phase 1 und 2)
Tageszeit / Ermüdungseffekte (falls relevant)

Zusätzlich interessant

Time-on-task Verteilung: Nicht nur Mittelwert, sondern auch Varianz — schreiben Journalisten mit KI gleichmäßiger?
Edit-Rate: Wie viel wird nach KI-Generierung noch verändert? Das sagst du zwar nicht direkt, aber wenn deine App es loggt wäre es Gold wert.


Perceived Text Quality
Hier wird es etwas kniffliger, weil "Qualität" subjektiv ist. Du brauchst eine klare Operationalisierung. Typische Ansätze:
Selbsteinschätzung der Journalisten (einfach, aber subjektiv)
Likert-Skala nach jeder Aufgabe, z.B.:

„Wie zufrieden sind Sie mit dem entstandenen Text?" (1–5)
„Wie sehr entspricht der Text Ihrem gewohnten Stil?" (1–5)

Externe Bewertung (aufwändiger, aber objektiver)
Redakteure oder andere Journalisten bewerten Texte blind — ohne zu wissen ob KI beteiligt war. Das wäre methodisch sehr stark.
Automatische Qualitätsmetriken (als Ergänzung)
Da du ohnehin NLP machst, kannst du das verbinden:

Kohärenzscores (z.B. mit einem Language Model)
Grammatikfehlerrate
\section{Writing Style Analysis}

\subsection{Lexical Features}
Lexikalische Features
Type-Token-Ratio (TTR) / MATTR
Misst lexikalische Vielfalt. MATTR (Moving Average TTR) ist robuster bei unterschiedlichen Textlängen — wichtig, weil Titel vs. Artikeltext riesige Längenunterschiede haben.
Lexikalische Dichte
Anteil der Inhaltswörter (Nomen, Verben, Adjektive, Adverbien) vs. Funktionswörter. KI-Texte tendieren oft zu höherer Dichte.
Nominalstil-Index
Verhältnis Nomen zu Verben. Journalismus hat traditionell hohen Nominalstil — verändert KI das?
Durchschnittliche Wortlänge & Wortfrequenz
Einfach zu berechnen, aber aussagekräftig. KI-Modelle bevorzugen statistische Häufigkeit, also tendenziell häufigere Wörter.
Hapax Legomena Rate
Anteil Wörter, die nur einmal vorkommen — Maß für Originalität/Kreativität.

Syntaktische Features
Durchschnittliche Satzlänge
Klassiker. KI-Texte sind oft gleichmäßiger — das ist selbst ein Signal.
Varianz der Satzlänge
Genauso wichtig wie der Mittelwert. Menschen variieren stärker.
Dependency-Tree-Tiefe
Mit spaCy oder CoreNLP einfach berechenbar. Misst syntaktische Komplexität.
Subordinationsgrad
Verhältnis von Neben- zu Hauptsätzen. Tiefere Schachtelung = komplexerer Stil.
POS-Tag-Distributionen
Anteil Adjektive, Adverbien, Passivkonstruktionen etc. als Feature-Vektor — gut für maschinelles Lernen.

Stilistische / höherwertige Features
Readability Scores
Flesch-Kincaid, Gunning Fog, oder für Deutsch: Lesbarkeitsindex (LIX) oder Wiener Sachtextformel. Gut für Vergleich, weil leicht erklärbar.
Stilhomogenisierung (dein Kern-Beitrag!)
Das ist dein stärkstes Argument: Schreiben verschiedene Journalisten ähnlicher, wenn sie KI nutzen? Du könntest paarweise Cosine-Similarity zwischen Autoren berechnen — wenn die Varianz zwischen Autoren in Phase 2 sinkt, ist das ein starkes Ergebnis.
Sentiment & Emotionalität
Mit Lexikon-basierten Tools (z.B. VADER, SentiWS für Deutsch) — verändert KI den emotionalen Ton?

Statistische Auswertung
Da du gepaarte Messungen hast (selber Journalist, Phase 1 vs. Phase 2), bieten sich an:

Wilcoxon-Vorzeichen-Rang-Test (non-parametrisch, kleine Stichprobe)
Paired t-Test wenn Normalverteilung annehmbar
Cohen's d für Effektgröße
PCA oder t-SNE um Texte im Feature-Raum zu visualisieren — sehr überzeugend für die Homogenisierungs-These


Tool-Empfehlungen
ToolSpracheFür wasspaCyPythonPOS-Tagging, Dependency ParsingNLTKPythonTTR, TokenisierungtextstatPythonReadability ScoresSentiWS / VADERPythonSentimentscikit-learnPythonCosine Similarity, PCAR (stylo package)RStilometrie, speziell für Autorenvergleich
\subsection{Syntactic Features}
\subsection{Readability}

\section{Qualitative Interview Findings}